
This paper tested whether community breadth at benchmark introduction --- operationalized as paper submission count at plurality-introduction year --- predicts benchmark displacement timing in Papers With Code. After validating all 5 FAIL FAST gates (G0 = 0.862, G1 = 0.605, G2 = 0.975, G3 = 0.246, G4 = 2.14), Cox regression on 258 complete-case rows (176 displacement events) yields:

> \textbf{HR = 1.006 (95\% CI = [0.846, 1.196]), LRT p = 0.9495}

This is a near-perfect null ($\Delta logL$ = 0.0020). Neither the lock-in nor the saturation mechanism is operative at a detectable level with this operationalization. Submission-count community breadth diversity is ruled out as a predictor class for benchmark displacement timing in the Papers With Code ecosystem (2015--2023). This result is specific to the submission-count operationalization; score velocity and institutional diversity remain untested at full power.

The 5-gate FAIL FAST protocol constitutes a reusable methodological contribution for benchmark lifecycle survival analysis, providing a pre-registration-equivalent framework for separating data quality failure from hypothesis failure.

Future directions: (1) score trajectory ($\Delta score$\_lag1\_z) as primary predictor on the full 258-row complete-case panel; (2) institution-deduplicated breadth via Semantic Scholar author API; (3) cross-platform replication on Semantic Scholar or OpenReview data.

